Cette section présente un résumé détaillé en français de ce tapuscrit de thèse. Ce résumé vise à rendre le contenu de la thèse plus accessible aux lecteurs francophones tout en contribuant à la valorisation de la langue française dans le contexte académique et scientifique. Dans les sections suivantes, nous présenterons l'objectif de cette thèse et soulignerons son importance, avant de proposer un résumé du contenu et des contributions de chaque chapitre.

\section*{Introduction}

Au cours des dernières décennies, les réseaux filaires et sans fil ont connu des avancées remarquables, révolutionnant la connectivité et permettant une large gamme d'applications dans divers secteurs. Les réseaux filaires ont évolué, passant des infrastructures à base de cuivre à la technologie avancée de fibre optique, offrant des vitesses de l'ordre du gigabit. Parallèlement, les réseaux sans fil ont également connu une croissance considérable, les réseaux mobiles évoluant de la 2G à la 5G, et les standards Wi-Fi progressant du 802.11b vers le Wi-Fi 6 et bientôt du Wi-Fi 7.

Ces avancées ont conduit à l'émergence des applications à faible latence, telles que les jeux en nuage ou jeux à la demande  (\textit{cloud gaming} en anglais), la réalité virtuelle et augmentée en nuage (\textit{cloud virtual reality} en anglais), ainsi que l'Internet tactile ou la chirurgie à distance. Ces applications sont considérées comme des innovations majeures qui transformeront les secteurs du divertissement, de la santé, de l’éducation et de l’industrie en offrant des expériences collaboratives en temps réel, des simulations immersives et des environnements virtuels interactifs.

Le succès de ces applications à faible latence repose sur des réseaux capables de fournir des vitesses élevées, une latence ultra-faible et une fiabilité accrue. Cependant, ces exigences posent d'importants défis dans des environnements réseau à capacité variable, comme les réseaux Wi-Fi ou les réseaux cellulaires comme la 4G ou la 5G. Les fluctuations de bande passante, la latence, le jitter ou encore les pertes de paquets sont des problèmes fréquents pouvant dégrader considérablement la qualité d’expérience des utilisateurs.

Ces défis rendent la détection des anomalies et le diagnostic des causes racines essentiels pour maintenir les performances des réseaux et garantir une qualité d'expérience satisfaisante. Les opérateurs réseau, tels qu’Orange, doivent disposer d’outils efficaces pour identifier rapidement les problèmes et les résoudre, surtout à l'ère des technologies immersives qui deviennent de plus en plus omniprésentes.


Cette thèse, menée dans le cadre d’un contrat CIFRE entre Orange Innovation et le centre INRIA de l'Université de Lorraine, vise à répondre aux limitations actuelles dans la détection d’anomalies et le diagnostic des causes racines dans les applications à faible latence. Réaliser cet objectif nécessite de résoudre quelques défis significatifs qui sont les suivants:

\begin{itemize}
    \item Collecter des datasets d'indicateurs de performances sur différentes applications à faible latence comme le cloud gaming ou le cloud en réalité virtuelle, dans des conditions réseaux réalistes et comportant des scénarios qui sont sujets à des dégradations de performance sur ces applications ;
    \item Développer des méthodes efficaces et robustes au phénomène de contamination de données. Ces méthodes permettraient de détecter les anomalies en utilisant des techniques d'apprentissage automatique non supervisées, à partir des différents indicateurs de performance collectés  ;
    \item Identifier les causes de dégradation de performance des applications à faible latence à l'aide de techniques de diagnostic reposant sur l'apprentissage automatique.
\end{itemize}

Pour répondre à ces problématiques, des contributions que nous détaillerons dans la suite, ont été effectuées:
\begin{itemize}
    \item La mise en place de bancs de test pour collecter des données pertinentes sur des plateformes de cloud gaming et de cloud VR opérant sous des réseaux à capacité variable ;
    \item L'évaluation de modèles d'apprentissage automatique non-supervisé pour la détection d'anomalies dans les métriques collectées sur notre banc de test de cloud gaming ;
    \item Le développement d'un algorithme à base d'apprentissage contrastif pour améliorer la détection d'anomalies et la robustesse à la contamination de données ;
    \item La conception d'un modèle à deux étapes pour le diagnostic des causes racines combinant apprentissage contrastif et classification supervisée pour identifier les sources spécifiques de dégradation lors de l'utilisation des applications de Cloud VR avec des réseaux Wi-Fi.
\end{itemize}

\section*{Etat de l'art}
Dans le Chapitre \ref{chap:related_work}, nous explorons les fondements et l'état de l'art des thématiques essentielles à cette thèse.

Tout d'abord nous nous sommes intéressés aux applications à faible latence. Ces applications  partagent un besoin commun : une latence de bout en bout minimale et des débits élevés. Dans le cas du cloud gaming et du cloud VR, qui sont au cœur de cette thèse, ces applications nécessitent une fiabilité réseau élevée pour fournir des expériences interactives et immersives. Leur fonctionnement repose sur des architectures et des protocoles complexes, où des facteurs comme le délai, la gigue, le débit et la perte de paquets jouent un rôle critique. De nombreuses études mettent en évidence l'impact de ces paramètres réseau sur la qualité d'expérience utilisateur.

Ensuite, nous avons présenté les avancées dans les réseaux à capacité variable, à savoir les réseaux cellulaires et Wi-Fi. L'évolution des réseaux cellulaires, de la 1G à la 5G, a permis des avancées significatives. La 5G se distingue par une latence ultra-faible, une bande passante accrue et des concepts innovants tels que le \textit{network slicing} ou le calcul en périphérie du réseau (\textit{edge computing}). Malgré ces avancées, les réseaux cellulaires continuent de faire face à des défis comme l’instabilité de la latence et l’indisponibilité de la bande passante, qui affectent les performances des applications à faible latence. Le même constat a été effectué avec les réseaux Wi-Fi où les normes, de l’IEEE 802.11 au Wi-Fi 6E, ont considérablement amélioré le débit et réduit la latence. Cependant, des limitations subsistent, notamment des interférences, des congestions, et des problématiques liées aux environnements denses qui contribuent à dégrader l'expérience utilisateur.

Les techniques d'apprentissage automatique ou profond, que nous utilisons pour répondre aux problématiques de cette thèse, et qui constituent une des thématiques clés sont présentées. Ces techniques jouent un rôle clé dans la résolution de problèmes complexes dans divers domaines, y compris les réseaux. L'apprentissage profond, basé sur les réseaux de neurones, est particulièrement performant pour capturer des motifs complexes dans les données. Il est utilisé selon plusieurs paradigmes d’apprentissage : supervisé, non supervisé, ou auto-supervisé. La conception d'une solution d'apprentissage profond implique plusieurs étapes à savoir la préparation des données, la conception du modèle, l'entraînement, l'évaluation et le déploiement qui garantissent la performance et la généralisation nécessaires pour des scénarios réels.

Nous nous sommes ensuite intéressés à la détection d'anomalies dans les séries temporelles qui vise à identifier les déviations par rapport aux schémas normaux, et qui est ainsi cruciale pour les systèmes de surveillance, notamment les réseaux. La littérature montre que ce problème est majoritairement abordés par des approches d'apprentissage non supervisées telles que les modèles statistiques, les techniques de regroupement (\textit{clustering}) ou des approches basées sur la reconstruction ou la prévision. Des progrès récents à base d’apprentissage auto-supervisé ou d’apprentissage contrastif, ont amélioré les capacités de détection d'anomalies. L’évaluation de ces techniques repose sur des métriques comme le F1-score, le MCC ou l’AU-PR.

La dernière thématique clé de cette thèse est le diagnostic des causes racines. Il vise à identifier les facteurs spécifiques responsables des dégradations de performance sur les réseaux à capacité variable. Traditionnellement, le diagnostic reposait sur des méthodes basées sur l'expertise humaine. Cependant, ces approches s'avèrent aujourd'hui inefficaces en raison de l'énorme quantité de données à traiter et sont désormais remplacées par des solutions automatisées et pilotées par l’apprentissage automatique telles que le clustering, les méthodes statistiques, les modèles d'apprentissage supervisé classiques ou profond.


\section*{Collecte de données sur les applications à faible latence}
Dans le Chapitre \ref{chap:measurements}, nous présentons en détails les méthodologies et les expérimentations mises en place pour collecter des données essentielles à l’évaluation des modèles d’apprentissage automatique destinés à la détection d’anomalies et au diagnostic des causes racines dans les applications à faible latence. Ces données constituent la base des analyses et des modèles développés dans les chapitres suivants.

Tout d'abord un banc de test a été mis en place pour collecter des fichiers d'opportunités de transmission permettant d'effectuer des expérimentations reproductibles sur les réseaux à capacité variable. Nous avons ainsi collecter des fichiers permettant de reproduire les conditions 4G du réseau opérationnel mobile d'Orange en utilisant un outil de génération appelé \textit{Saturator}, sous différents niveau de qualité de connectivité 4G.

A partir de ces fichiers, un banc de test a été mis en place pour collecter des données sur trois (03) plateformes commerciales de cloud gaming (\textit{Stadia} de Google, \textit{GeForce Now} de NVIDIA et \textit{XCloud} de Microsoft) en conditions de réseau 4G émulées en utilisant l'outil \textit{Mahimahi-Linkshell} et l'outil open-source \textit{DECAF}. Nous avons ainsi collecté des séries temporelles multivariées contenant des métriques réseaux et applicatives reflétant les performances des plateformes de cloud gaming sous différents scénarios.

Le troisième et dernier banc de test développé était dédié aux applications de Cloud VR opérant sous des réseaux Wi-Fi. Ces expérimentations ont inclus à la fois des conditions normales et des scénarios de dégradation émulés en laboratoire, tels que des interférences et une atténuation du signal. Les métriques collectées comprennent des métriques applicatives issues des outils Oculus OVR Metrics et CloudXR de NVIDIA, incluant des indicateurs de performance de la session de Cloud VR, des métriques réseau capturées directement depuis le routeur Livebox, et des données réseaux brutes PCAP collectées à l’aide de sondes placées dans des cages de Faraday près des casques VR pour analyser les interactions au niveau des paquets entre les clients VR et les serveurs.

Ces jeux de données représentent des ressources précieuses à l'étude des performances des applications à faible latence.  Ils permettent de modéliser des conditions réseau réalistes, d’évaluer et de comparer les performances des modèles de détection d’anomalies et de diagnostic des causes racines et de tester la généralisabilité des solutions proposées dans cette thèse.

\section*{Evaluation de modèles non-supervisé pour la détection d'anomalies}

Le Chapitre \ref{chap:ml_evaluation} s'intéresse à l'application et l'évaluation de modèles d'apprentissage non-supervisé pour la détection d'anomalies dans les séries temporelles multivariées. L’objectif principal est de comparer différentes approches pour identifier les anomalies dans les indicateurs de performance des applications à faible latence, en mettant un accent particulier sur l'efficacité de détection de ces méthodes, leur robustesse à la contamination de données et leur efficacité en temps de calcul et d'inférence.

Plusieurs modèles d'apprentissage non-supervisé existent dans la littérature et ont été evalués une sélection des plus utilisés, notamment les forêts d'isolation (\textit{Isolation Forest}), les algorithmes de classification à une classe (\textit{OC-SVM, Deep-SVDD}), les approches de types reconstruction à base d'analyse en composantes principales (\textit{PCA}) ou d'autoencodeurs(\textit{AE, VAE, DAGMM, USAD}). Afin d'effectuer une évaluation juste et cohérente de ces algorithmes, une méthode nommée \textit{\ac{WAD}} a été introduite pour améliorer l'évaluation et la performance des algorithmes a été mesurée avec des métrique comme le F1-score et le \ac{MCC}.

Cette évaluation a ainsi montré que les modèles comme l'Isolation Forest, DAGMM et USAD présentent de bonnes performances pour détecter des anomalies. Aussi, nous avons noté que les modèles voient leur performance se dégrader en présence de données d'entraînement contaminées. Un autre résultat notable est que le modèle le plus rapide à l'entraînement et à l'inférence est le PCA, qui n'a malheureusement pas d'aussi bonnes performances en détection d'anomalies que DAGMM et USAD, qui eux ont une efficacité de calcul modérée.

Ce chapitre est conclu par des recommandations pratiques sur l’utilisation des modèles d’apprentissage non supervisé pour la détection d’anomalies dans des scénarios réels.


% \section*{Détection du trafic cloud gaming par apprentissage non-supervisé}

% Le Chapitre \ref{chap:cg_detection} explore la détection du trafic cloud gaming dans les environnements réseau, une tâche essentielle pour optimiser la latence via des méthodes d'ingénierie de trafic et améliorer la qualité d'expérience des utilisateurs. Contrairement aux approches traditionnelles basées sur des méthodes supervisées, ce travail redéfinit la détection du trafic CG comme un problème d’apprentissage non supervisé, permettant une meilleure généralisation à des flux inconnus. 

% En utilisant les résultats du chapitre précédent \ref{chap:ml_evaluation}, nous avons sélectionné le modèle non-supervisé USAD qui identifie les flux cloud gaming en exploitant uniquement les caractéristiques intrinsèques des données. Ce modèle de détection a été intégré dans une architecture hybride en vue de son déploiement. Cette architecture comprend un module matériel (routeur P4) qui effectue un traitement rapide des paquets. Il extrait les métriques utilisés par le modèle USAD pour la détection et les transmets au module de virtualisation des fonctions réseau dans lequel tourne le modèle de détection pour classer les flux à partir de ces métriques.

% Cette approche a été évalué à partir de métriques extraites de traces réseau issues de plateformes commerciales de cloud gaming. L'évaluation montre que notre approche surpasse l'approche supervisée (arbre de décision) lorsqu'ils sont confrontés à des traces issues de nouvelles plateformes de cloud gaming ou issues de nouveaux jeux. L'architecture hybride permet de coupler vitesse de traitement grâce au module P4 et précision de la détection.


\section*{Une nouvelle approche de détection d'anomalies par apprentissage contrastif}

Le Chapitre \ref{chap:cats} présente CATS (\textit{Contrastive Anomaly detection on Time Series}), une approche non-supervisée d'apprentissage conçu pour améliorer la détection d’anomalies dans les séries temporelles multivariées. CATS est une contribution majeure de cette thèse, apportant une approche innovante basée sur l’apprentissage contrastif pour surmonter les limitations des méthodes existantes en termes de robustesse à la contamination des données, et de performance sur les séries temporelles. CATS repose sur plusieurs innovations clés :
\begin{itemize}
    \item L'introduction d'une nouvelle fonction de perte contrastive utilisant la similarité temporelles basée sur la métrique temporelle DTW (\textit{Dynamic Time Warping}). Cette fonction de perte contribue à regrouper dans une portion de l'espace latent les fenêtres temporelles similaires d'un point de vue temporel tout en éloignant les autres ;
    \item La génération d'anomalies synthétiques dans le jeu d'entraînement à l’aide de techniques de transformation de données telles que le masquage et le bruitage, renforçant ainsi la robustesse du modèle face aux données contaminées ;
    \item CATS est agnostique à l’architecture neuronale utilisée pour encoder les séries temporelles. Il peut être combiné avec des architectures comme des autoencodeurs, des réseaux de neurones récurrents, ou des Transformers pour capturer des dépendances locales et globales.
\end{itemize}

Evalué en utilisant des jeux de données publics et nos données collectées précédemment pour le cloud gaming, CATS montre qu'il surpasse les modèles de l'état de l'art en terme de précision et aussi de robustesse à la contamination de données. La principale limite de CATS reste son coût computationnel élevé en raison de la métrique DTW qui limite son applicabilité et aussi l'étendue de son efficacité.

\section*{Diagnostic des causes racines dans le cloud en réalité virtuelle}

Le Chapitre \ref{chap:diagnosis} se concentre sur le diagnostic des causes racines dans les applications de réalité virtuelle en cloud fonctionnant sur des réseaux Wi-Fi. Il propose un algorithme à deux étages pour détecter et diagnostiquer les dégradations de performance, en combinant la détection d'anomalies avec un classifieur supervisé. Cette approche vise à relever les défis liés à l’identification des causes spécifiques des dégradations dans des environnements réseau dynamiques.

Les deux étages de cette approche sont les suivantes:
\begin{itemize}
    \item La détection d'anomalies en utilisant le modèle CATS (présenté au chapitre précédent) qui donnera un score d'anomalies qui servira à l'étape suivante.
    \item La classification des causes à l'aide d'un classifier supervisé tel qu'une machine à vecteurs de supports (\textit{Support Vector Machine}) pour mapper les anomalies précédemment détectées à des classes spécifiques de causes (par exemple les interférence ou l'atténuation).
\end{itemize}

Les expérimentations ont été réalisées en utilisant les données collectées dans le Chapitre \ref{chap:measurements} sur le cloud VR, notamment les métriques réseau capturées au niveau de la Livebox. Les résultats montrent que notre approche surpasse les approches de classification de la littérature dans le diagnostic des causes de dégradation dans le cloud VR. Elle démontre aussi de bonnes performances même dans les scénarios ou peu de données étiquetées sont disponibles grâce à l'efficacité du détecteur d'anomalies. Cette approche modulaire présente aussi une faible complexité computationnelle la rendant applicable à des déploiements en temps réel.


\section*{Conclusion et perspectives}

Cette thèse a abordé le problème de la détection d'anomalies et le diagnostic des applications à faible latence. Cependant, de nouvelles opportunités de recherche se dessinent pour approfondir et élargir les contributions présentées. Tout d'abord, des collectes de données supplémentaires pourraient être réalisées dans des environnements Cloud VR, en explorant des scénarios de dégradation de la qualité Wi-Fi plus complexes ou en effectuant des collectes à large échelle dans des réseaux domestiques ou d’entreprise. Ensuite, l’application de notre solution de diagnostic proposée à des réseaux 5G, où les dégradations sont spécifiques, permettrait d’évaluer leur généralisation dans d'autres environnements.

Par ailleurs, l’intégration de multiples sources de données (par exemple, les données PCAP, les métriques client et celles des points d'accès) grâce à des approches d’apprentissage multi-modales pourrait améliorer le diagnostic. Enfin, des techniques comme la découverte de classes inconnues et la découverte causale pourraient être explorées pour détecter des causes de dégradations encore non identifiées et comprendre les relations causales entre les métriques réseau, renforçant ainsi l’explicabilité des solutions proposées. 
